% !TEX root = ../main.tex
% ---------------------------------------------------------------
% GENERATED FILE - DO NOT EDIT.
% Produced by docs/tools/render_reference.py from the repository source.
% Regenerate with: make docs
% ---------------------------------------------------------------

\apibreak
\section{\texttt{validate\_feedback\_csv}}
\label{func:feedback_intelligence_agent.data_contracts.validate_feedback_csv}\index{Functions!validate\_feedback\_csv}

Validate a feedback CSV file against the data contract.
\apifield{Usage}
\begin{lstlisting}[style=usage, language=Python]
validate_feedback_csv(path: str | Path, *, strict: bool = False) -> tuple[ValidationReport, list[FeedbackRecord]]
\end{lstlisting}
\apifield{Arguments}
\begin{arglist}
  \item[\texttt{path}] \texttt{str \textbar{} Path}, required. Path to the CSV file to validate.
  \item[\texttt{strict}] \texttt{bool}, default \texttt{False}, keyword-only. When True, raise \texttt{DataContractError} if any error is found. When False, invalid rows are skipped and reported.
\end{arglist}
\apifield{Value}\texttt{tuple[ValidationReport,\allowbreak{} list[FeedbackRecord]\allowbreak{}]\allowbreak{}}. A tuple of the validation report and the feedback records that passed validation (all rows in a clean dataset, the valid subset otherwise).
\apifield{Raises}
\begin{arglist}
  \item[\texttt{FileNotFoundError}] If the file does not exist.
  \item[\texttt{DataContractError}] In strict mode, if the dataset has any validation error.
  \item[\texttt{FileReadError}] If an existing file cannot be read.
  \item[\texttt{DataLoadingError}] If the CSV cannot be parsed.
\end{arglist}
\apifield{Source}\texttt{src/\allowbreak{}feedback\_\allowbreak{}intelligence\_\allowbreak{}agent/\allowbreak{}data\_\allowbreak{}contracts.\allowbreak{}py:120}~(\href{https://github.com/DiogoRibeiro7/feedback-intelligence-agent/blob/edc9d78403eeed9f7d6ee5aef6b8d9b1b5780ef4/src/feedback_intelligence_agent/data_contracts.py#L120}{online}), module \cref{module:feedback_intelligence_agent.data_contracts}
